% !TeX program = pdflatex
\documentclass[preprint,12pt]{elsarticle}

% -------------------------------------------------
% Packages
% -------------------------------------------------
\usepackage[T1]{fontenc}
\usepackage[utf8]{inputenc}
\usepackage{amsmath,amssymb}
\usepackage{graphicx}
\usepackage{booktabs}
\usepackage{hyperref}
\usepackage{microtype}
\renewcommand{\ttdefault}{cmtt}
\microtypesetup{expansion=false}

\hypersetup{
  colorlinks=true,
  linkcolor=blue,
  citecolor=blue,
  urlcolor=blue
}

\renewcommand{\thefigure}{S\arabic{figure}}
\renewcommand{\thetable}{S\arabic{table}}
\renewcommand{\thesection}{S\arabic{section}}
\newcommand{\suppwideimage}[1]{\includegraphics[width=\linewidth,height=0.82\textheight,keepaspectratio]{#1}}

\begin{document}

\begin{frontmatter}

\title{Supplementary Material for\\
Secretary-Style Split Search for Decision Trees}

\author{Mathias Valla}

\begin{abstract}
This supplementary document provides the additional figures supporting the main manuscript: Supplementary Figure~S1 on size-stratified Pareto comparisons, Supplementary Figure~S2 on the screening analysis of the parametric secretary strategy, Supplementary Figures~S3--S5 on ridgeline summaries of within-dataset variability, Supplementary Figures~S6--S7 on reliability and joint-success plots, Supplementary Figure~S8 on tolerance-constrained quickest-method regime maps, Supplementary Figures~S9--S11 on within-dataset whisker summaries, Supplementary Figure~S12 on exhaustive gain-landscape diagnostics for the regression--classification asymmetry, and Supplementary Figures~S13--S14 on split-search effort comparisons using the selected \(S_{\mathrm{par}}\) representative.
\end{abstract}

\end{frontmatter}

\begin{figure}[p]
\centering
\suppwideimage{SUPP_FIGURES/supp_figure_01_pareto_large_small.png}
\caption{Size-stratified Pareto comparisons. The top row reports large datasets, defined by \(n\times p \ge 25000\), and the bottom row reports small datasets, defined by \(n\times p \le 2500\). Columns correspond to regression, classification with Gini impurity, and classification with entropy impurity. In each panel, each faint point represents one benchmark dataset summarized by the median over 100 runs for a given method variant. The horizontal axis is the median percentage of training time saved relative to the exhaustive \texttt{scikit-learn} splitter, and the vertical axis is the median percentage loss relative to the same baseline. The black point at the origin is the exhaustive splitter. Large white-edged markers indicate the cross-dataset centroids. Faint Pareto frontiers summarize the centroid-level and dataset-level envelopes.}
\label{fig:supp_pareto_size}
\end{figure}

\newpage

\begin{figure}[p]
\centering
\suppwideimage{SUPP_FIGURES/supp_figure_02_secretary_par_all_large_small.png}
\caption{Screening analysis of the parametric secretary strategy \(S_{\mathrm{par}}\). Rows correspond to all datasets, large datasets, and small datasets. Columns correspond to regression, classification with Gini impurity, and classification with entropy impurity. Each colored point represents one tested configuration of \(S_{\mathrm{par}}\), defined by the number of gain samples used to fit the working distribution and by the acceptance quantile used for stopping. The horizontal axis reports the median percentage of training time saved relative to the exhaustive splitter, and the vertical axis reports the median percentage loss relative to the same baseline. The black point at the origin is the exhaustive baseline, and large white-edged markers indicate cross-dataset centroids of the corresponding configurations.}
\label{fig:supp_spar_screening}
\end{figure}

\newpage

\begin{figure}[p]
\centering
\suppwideimage{SUPP_FIGURES/supp_figure_03_ridgelines_regression.png}
\caption{Ridgeline summaries for regression. Each horizontal ridge corresponds to one benchmark dataset, identified on the vertical axis by its sample size and number of features. The left panel shows the empirical distribution, over 100 runs, of the percentage of training time saved relative to the exhaustive splitter; the right panel shows the corresponding distribution of percentage RMSE loss relative to the same baseline. Colors identify methods, and the dashed vertical line at zero marks parity with the exhaustive splitter.}
\label{fig:supp_ridgeline_regression}
\end{figure}

\newpage

\begin{figure}[p]
\centering
\suppwideimage{SUPP_FIGURES/supp_figure_04_ridgelines_classification_gini.png}
\caption{Ridgeline summaries for classification with Gini impurity. Each horizontal ridge corresponds to one benchmark dataset. The left panel shows the empirical distribution, over 100 runs, of the percentage of training time saved relative to the exhaustive splitter, and the right panel shows the corresponding distribution of percentage F1 loss relative to the same baseline. Colors identify methods, and the dashed vertical line at zero marks parity with the exhaustive splitter.}
\label{fig:supp_ridgeline_gini}
\end{figure}

\newpage

\begin{figure}[p]
\centering
\suppwideimage{SUPP_FIGURES/supp_figure_05_ridgelines_classification_entropy.png}
\caption{Ridgeline summaries for classification with entropy impurity. As in Supplementary Figure~\ref{fig:supp_ridgeline_gini}, each horizontal ridge corresponds to one benchmark dataset, the left panel shows percentage time saved, and the right panel shows percentage F1 loss, both relative to the exhaustive splitter and computed over 100 repeated runs.}
\label{fig:supp_ridgeline_entropy}
\end{figure}

% \newpage

% \begin{figure}[p]
% \centering
% \includegraphics[width=0.6\linewidth]{SUPP_FIGURES/supp_figure_06_success_joint_large_small.png}
% \caption{Joint success probabilities stratified by dataset size. The upper half of the figure reports large datasets and the lower half reports small datasets. Columns correspond to predictive-loss tolerances \(\varepsilon \in \{5\%,10\%,20\%\}\), and within each size stratum the rows correspond to minimum time-saving thresholds \(s \in \{0\%,10\%,25\%,50\%\}\). In each panel, bar heights give the empirical fraction of runs satisfying the joint event ``time saved \(\ge s\) and predictive loss \(\le \varepsilon\)'', aggregated separately for regression, Gini classification, and entropy classification. Colors identify methods. These summaries complement the continuous success curves by highlighting practically relevant operating points.}
% \label{fig:supp_joint_success_size}
% \end{figure}

\newpage

\begin{figure}[p]
\centering
\suppwideimage{SUPP_FIGURES/supp_figure_07_success_loss_only_large_small.png}
\caption{Loss-only success probabilities stratified by dataset size. The top row reports large datasets and the bottom row reports small datasets; columns correspond to regression, Gini classification, and entropy classification. In each panel, the horizontal axis is the predictive-loss tolerance \(\varepsilon\), and the vertical axis is the empirical fraction of runs whose loss is at most \(\varepsilon\), aggregated over the datasets in the corresponding size stratum. Colored curves identify methods, and the horizontal black line at one is the exhaustive baseline. All curves are computed over 100 repeated runs.}
\label{fig:supp_loss_success_size}
\end{figure}

\begin{figure}[p]
\centering
\suppwideimage{SUPP_FIGURES/supp_figure_12_success_joint_all.png}
\caption{Joint success probabilities over all benchmark datasets. Columns correspond to predictive-loss tolerances \(\varepsilon \in \{0.5\%,1\%,2.5\%\}\), and rows correspond to minimum time-saving thresholds \(s \in \{0\%,10\%,25\%,50\%\}\). In each panel, bar heights report the empirical fraction of runs satisfying the joint event ``time saved \(\ge s\) and predictive loss \(\le \varepsilon\)'', aggregated over all datasets for regression, classification with Gini impurity, and classification with entropy impurity. All probabilities are estimated over 100 repeated runs.}
\label{fig:supp_joint_success_all}
\end{figure}

\newpage

\begin{figure}[p]
\centering
\suppwideimage{SUPP_FIGURES/supp_figure_08_predicted_regime_maps.png}
\caption{Tolerance-constrained quickest-method regime maps in the \((\log n, p/n)\) plane. Rows correspond to regression, Gini classification, and entropy classification, and columns correspond to predictive-loss tolerances \(\tau \in \{0.5\%,1\%,2.5\%\}\). At each dataset location, the color marks the fastest method whose median loss stays at or below the displayed tolerance; when no non-exhaustive method satisfies the tolerance, the exhaustive splitter is shown instead.}
\label{fig:supp_regime_maps}
\end{figure}

\newpage

\begin{figure}[p]
\centering
\suppwideimage{SUPP_FIGURES/supp_figure_09_within_whiskers_loss_speedup.png}
\caption{Within-dataset whisker summaries for regression. The upper half of the figure reports dataset-level median speedups for selected methods, and the lower half reports the corresponding dataset-level median RMSE losses. Each point represents one benchmark dataset, ordered within each method panel. Horizontal whiskers indicate \(\pm 0.5\) times the empirical interquartile range over the 100 repeated runs.}
\label{fig:supp_within_regression}
\end{figure}

\newpage

\begin{figure}[p]
\centering
\suppwideimage{SUPP_FIGURES/supp_figure_10_within_whiskers_loss_speedup.png}
\caption{Within-dataset whisker summaries for classification with Gini impurity. As in Supplementary Figure~\ref{fig:supp_within_regression}, the upper half reports dataset-level median speedups and the lower half reports dataset-level median F1 losses for a selected subset of methods. Each point is one dataset and the whiskers indicate \(\pm 0.5\) times the empirical interquartile range across the 100 repeated runs.}
\label{fig:supp_within_gini}
\end{figure}

\newpage

\begin{figure}[p]
\centering
\suppwideimage{SUPP_FIGURES/supp_figure_11_within_whiskers_loss_speedup.png}
\caption{Within-dataset whisker summaries for classification with entropy impurity. The upper half reports dataset-level median speedups and the lower half reports dataset-level median F1 losses, with points corresponding to datasets and whiskers indicating \(\pm 0.5\) times the empirical interquartile range over the 100 repeated runs.}
\label{fig:supp_within_entropy}
\end{figure}

\newpage

\begin{figure}[p]
\centering
\suppwideimage{SUPP_FIGURES/supp_figure_13_gain_landscape.png}
\caption{Exhaustive gain-landscape diagnostics by task. Each box summarizes one benchmark task using dataset-level medians of internal-node diagnostics computed under exhaustive split search. The left panel reports the proportion of admissible thresholds whose gain lies within 5\% or 10\% of the nodewise optimum, the middle panel reports the relative gap between the best gain and the median gain of a random admissible threshold, and the right panel reports the contiguous width of the near-optimal region on the winning covariate.}
\label{fig:supp_gain_landscape}
\end{figure}

\newpage

\begin{figure}[p]
\centering
\suppwideimage{SUPP_FIGURES/supp_figure_14_effort_loss.png}
\caption{Effort--loss comparison by task. Each faint point represents one benchmark dataset summarized by the median over 100 runs for a given method, and each large white-edged marker represents the corresponding cross-dataset centroid. The horizontal axis reports the median percentage of gain evaluations saved relative to the exhaustive splitter, and the vertical axis reports the corresponding median predictive loss. Faint dashed frontiers summarize the dataset-level Pareto envelope, and darker solid frontiers summarize the centroid-level envelope. Only one \(S_{\mathrm{par}}\) representative, \(S_{\mathrm{par}}(\sqrt{n},q=0.75)\), is shown.}
\label{fig:supp_effort_loss}
\end{figure}

\newpage

\begin{figure}[p]
\centering
\suppwideimage{SUPP_FIGURES/supp_figure_15_time_effort_calibration.png}
\caption{Calibration of wall-clock savings against split-search effort. Each faint point represents one benchmark dataset summarized by the median over 100 runs for a given method, and each large white-edged marker represents the corresponding cross-dataset centroid. The top row compares median percentage time saved with median percentage total effort saved, where total effort is measured through the number of gain evaluations per tree. The bottom row compares median percentage time saved with median percentage effort saved per split call. Faint dashed frontiers summarize the dataset-level Pareto envelope, and darker solid frontiers summarize the centroid-level envelope. Only one \(S_{\mathrm{par}}\) representative, \(S_{\mathrm{par}}(\sqrt{n},q=0.75)\), is shown.}
\label{fig:supp_time_effort}
\end{figure}

\end{document}
